\section{The actual translated-initial comparison in the original unshifted coordinate ring}

\begin{theorem}[The whole unshifted original Taylor extension preserves every intrinsic Hilbert--Samuel value]
\label{hir:actual-unshifted-original-formal-hs-transfer}
\uses{hir:actual-original-formal-entire-hs-transfer}
\lean{Hironaka.HironakaActualUnshiftedInitialHSTransfer.hilbertSamuelValue_eq_actual_unshifted_formal_extension}
\leanok
For every ideal $I$ of the original analytic germ ring on the actual
$n$-dimensional model and every order $\ell$, its intrinsic Hilbert--Samuel
value equals that of the entire formal extension under the same actual
unshifted Taylor homomorphism $T$.
\end{theorem}
\begin{proof}
The proved original whole-quotient Taylor algebra equivalence gives this
identity on the unchanged original model of dimension $0+n$. Transport
that theorem and the actual original Taylor homomorphism along the true
dimension equality $0+n=n$. This is exactly the actual unshifted Taylor
homomorphism already defined on the original model. Both whole ideals
and their actual complex quotient actions are transported together, so
the quotient dimensions and the entire extension in the statement are
the same original quantities.
\end{proof}

\begin{definition}[Actual translated finite initial ideals in the original coordinate ring]
\label{hir:actual-unshifted-translated-initial-ideal}
\lean{Hironaka.HironakaActualUnshiftedInitialHSTransfer.actualUnshiftedTranslatedPolynomialEval,Hironaka.HironakaActualUnshiftedInitialHSTransfer.actualUnshiftedTranslatedInitialIdeal}
\leanok
For the true vector $p$, evaluate a finite complex polynomial at the
original coordinate germs $x_i+p_i$. For an actual formal ideal $K$,
let its translated initial ideal be the span of these evaluations of the
unchanged finite actual initial vertex polynomials of $K$. The whole
original analytic ideal is independent of the Hilbert--Samuel order.
\end{definition}

\begin{theorem}[The entire actual translated analytic ideal has its unchanged formal extension]
\label{hir:actual-unshifted-translated-initial-formal-extension}
\uses{hir:actual-unshifted-translated-initial-ideal}
\lean{Hironaka.HironakaActualUnshiftedInitialHSTransfer.actualUnshiftedTranslatedInitialIdeal_formal_extension}
\leanok
The whole Taylor extension of this actual analytic ideal is precisely
the existing actual translated initial formal ideal of $K$ at $p$.
\end{theorem}
\begin{proof}
The genuine Taylor coordinate dictionary sends the original coordinate
germs to the same actual formal variables, and its complex algebra map
preserves the constants $p_i$. Polynomial algebra extensionality therefore
identifies both entire finite polynomial evaluation maps. Ideal image
commutes with the span of their same finite vertex family. This proves
both directions of the entire ideal equality.
\end{proof}

\begin{theorem}[All translated-initial values are bounded by the unchanged actual original values]
\label{hir:actual-unshifted-translated-initial-all-hs}
\uses{hir:actual-unshifted-original-formal-hs-transfer,hir:actual-unshifted-translated-initial-formal-extension,hir:actual-translated-initial-all-order-hs-dominance}
\lean{Hironaka.HironakaActualUnshiftedInitialHSTransfer.hilbertSamuelValue_actual_unshifted_translated_initial_le,Hironaka.HironakaActualUnshiftedInitialHSTransfer.hilbertSamuelValue_actual_original_translated_initial_le}
\leanok
For every actual formal ideal $K$, every $p$ and every $\ell$, the actual
translated analytic initial ideal has value at most the genuine formal
value of $K$. In particular, for $K=T(I)$, it has value at most the
unchanged intrinsic original value $H_I(\ell)$, for every original
analytic ideal $I$.
\end{theorem}
\begin{proof}
Use the first theorem to identify the analytic value with that of its
whole actual formal extension. The preceding entire ideal dictionary
identifies this extension with the existing translated formal initial
ideal. The already proved genuine translated-initial comparison bounds
its value by that of $K$. In the case $K=T(I)$, apply the first theorem
again to identify the latter value with the original analytic value.
No comparison, coefficient convergence, or ideal-identification premise
is supplied by the caller.
\end{proof}
